\documentclass[10pt]{article}
\usepackage[margin=1in]{geometry}
\usepackage{times}
\usepackage{amsmath,amssymb}
\usepackage{graphicx}
\usepackage{booktabs}
\usepackage{xcolor}
\usepackage{subcaption}
\usepackage[colorlinks=true,allcolors=blue]{hyperref}
\usepackage{authblk}
\usepackage{natbib}

% placeholder macros — overwritten by make_paper_results.py
\newcommand{\brcaNmis}{2086}
\newcommand{\brcaRING}{642}
\newcommand{\brcaBRCT}{1363}
\newcommand{\evoModel}{\texttt{evo2\_7b}}
\newcommand{\esmModel}{\texttt{esm2\_t33\_650M}}
\newcommand{\dnaLayer}{\texttt{blocks.24.mlp.l3}}
\newcommand{\protLayer}{\texttt{layer 30}}
\newcommand{\Kshared}{64}
\newcommand{\Kprivate}{64}
\newcommand{\topk}{16}
\newcommand{\retRone}{--}
\newcommand{\retRoneCCA}{--}
\newcommand{\npca}{128}
 % auto-generated numeric macros

\title{\textbf{CentralDogma-$\Delta$CC: Disentangling Shared and Modality-Specific\\
Variant Mechanisms across Genomic and Protein Language Models}}
\author{Anonymous}
\affil{}
\date{}

\begin{document}
\maketitle

\begin{abstract}
Genomic language models (e.g.\ Evo\,2) and protein language models (e.g.\ ESM-2) both
predict variant effects, but it is unknown whether they encode the \emph{same} molecular
mechanisms internally, or complementary modality-specific ones. We introduce
\textbf{CentralDogma-$\Delta$CC}, a sparse \emph{cross-architecture, cross-modality
crosscoder} that operates not on raw activations but on \emph{paired variant-induced
activation deltas}: for each missense variant we compute
$\Delta h_{\mathrm{DNA}}=h_{\mathrm{DNA}}(\text{mut})-h_{\mathrm{DNA}}(\text{wt})$ in Evo\,2 and
$\Delta h_{\mathrm{PROT}}$ in ESM-2, and learn shared and modality-private sparse latents
that reconstruct both. On the BRCA1 saturation-genome-editing benchmark
(\brcaNmis{} missense variants with deep mutational scanning function scores) the shared
code (i)~aligns the two models well enough to retrieve the paired variant across modalities
far above chance (Recall@1 \retRone{} vs.\ CCA \retRoneCCA{}), (ii)~predicts DMS function
scores under a residue-disjoint split competitively with a concatenation baseline while
using \Kshared{} interpretable dimensions, and (iii)~yields latents that are enriched for
protein domains and loss-of-function. Causally ablating a shared latent's decoder direction
inside \emph{both} models moves their variant scores together, evidence that the latent
captures a genuinely shared mechanism rather than a correlational artifact. We release code,
activations, and the full evaluation harness.
\end{abstract}

\section{Introduction}
Foundation models now span the central dogma: Evo\,2 \citep{brixi2026evo2} models DNA at
single-nucleotide resolution, while ESM-2 \citep{lin2023esm2} models proteins. Both assign
scores to human coding variants, and both are used for variant effect prediction (VEP).
Yet a basic mechanistic question is open: when a DNA model and a protein model both flag a
missense variant, are they responding to the \emph{same} underlying cause (e.g.\ disruption
of a conserved structural residue), or to different, modality-specific causes (splicing and
codon context on the DNA side; folding and solvent exposure on the protein side)? Answering
this requires comparing the models' \emph{internal representations}, not just their output
scores. This is hard because the two models differ in architecture (StripedHyena vs.\
Transformer), tokenization (nucleotides vs.\ amino acids), input length, hidden width, and
training objective; naively aligning raw activations is dominated by global nuisance factors
such as gene identity and conservation.

We make three observations. First, the object of interest is not the absolute
representation of a sequence but the \emph{change} a variant induces; taking the paired
wild-type$\to$mutant activation delta in each model cancels most sequence-identity nuisance
and focuses on the variant's mechanism. Second, sparse dictionary methods (SAEs,
crosscoders) provide an interpretable, causally manipulable basis; crosscoders in particular
were designed to \emph{diff} representations across models \citep{lindsey2024crosscoders}.
Third, restricting to a single gene (BRCA1) removes the dominant gene-identity shortcut,
so any cross-modal alignment the model finds must reflect within-gene variant structure.

Combining these, \textbf{CentralDogma-$\Delta$CC} is a crosscoder over paired deltas with
an explicit \emph{shared / DNA-private / protein-private} factorization. Our contributions:
\begin{itemize}\itemsep2pt
\item A crosscoder applied across \emph{heterogeneous biological foundation models}
(a DNA and a protein model), rather than across layers or across a base/fine-tune pair.
\item Learning on \emph{paired variant-induced deltas} with an explicit shared--private split,
trained with reconstruction, cross-modal alignment, contrastive, and orthogonality losses
under BatchTopK sparsity.
\item A rigorous BRCA1 evaluation---cross-modal retrieval, residue- and domain-disjoint DMS
and ClinVar prediction, annotation enrichment, and a \emph{cross-model causal ablation} that
perturbs the same latent inside both Evo\,2 and ESM-2.
\end{itemize}

\section{Related Work}
\textbf{Genomic and protein language models.} Evo\,2 \citep{brixi2026evo2} is a StripedHyena
genomic model with strong zero-shot VEP, including on BRCA1 SGE data; ESM-2
\citep{lin2023esm2} is a protein masked-LM whose representations encode structure and
function. \textbf{Sparse features and crosscoders.} Sparse autoencoders recover
interpretable features in LMs and, via InterPLM \citep{simon2025interplm}, in protein LMs.
Crosscoders \citep{lindsey2024crosscoders} extend SAEs to jointly encode multiple
representation spaces and were introduced for cross-layer and model-diffing analysis; recent
work applies them to cross-architecture diffing \citep{jiralerspong2026crossarch} and to
narrow fine-tuning \citep{kassem2026delta,minder2025chat}. We instead diff across the
central dogma---a DNA and a protein model of the \emph{same} variant. \textbf{Variant effect
prediction.} DMS assays such as BRCA1 SGE \citep{findlay2018brca1} provide continuous
functional labels; we use them not as a leaderboard target but to test whether shared latents
carry mechanism.

\section{Method}
\label{sec:method}
\paragraph{Paired variant deltas.} For a missense SNV we build an 8{,}192\,bp
variant-centred window and score wild-type and mutant sequences with Evo\,2, extracting the
residual-stream activation at layer $\ell_{\mathrm{DNA}}$ and SNV position; the DNA delta is
$\Delta h_{\mathrm{DNA}}=h^{\text{mut}}-h^{\text{wt}}\in\mathbb{R}^{d_{\mathrm{DNA}}}$
($d_{\mathrm{DNA}}{=}4096$). Independently we substitute the residue in the protein and take
the ESM-2 hidden-state delta $\Delta h_{\mathrm{PROT}}\in\mathbb{R}^{d_{\mathrm{PROT}}}$
($d_{\mathrm{PROT}}{=}1280$) at the variant residue and layer $\ell_{\mathrm{PROT}}$. Deltas
are standardized per dimension and projected onto the top \npca{} principal components fit on
the training split (variant deltas are high intrinsic dimension, so this denoising is needed
for the reconstruction objective to generalize); decoder directions are mapped back to raw
activation units for the causal interventions.

\paragraph{Shared--private crosscoder.} Each modality $m\in\{\mathrm{DNA},\mathrm{PROT}\}$ has a
shared encoder/decoder and a private encoder/decoder:
\[
z^{m}_{s}=\mathrm{ReLU}(E^{m}_{s}(\Delta h_m-b_m)),\quad
z^{m}_{p}=\mathrm{ReLU}(E^{m}_{p}(\Delta h_m-b_m)),\quad
\widehat{\Delta h}_m = D^{m}_{s}\,\tilde z^{m}_{s} + D^{m}_{p}\,\tilde z^{m}_{p} + b_m,
\]
where $\tilde z$ applies BatchTopK \citep{bussmann2024batchtopk} (keep the $k\!\cdot\!B$
largest activations per code block across the batch). Shared codes of the same variant are
aligned across modalities. The training objective is
\[
\mathcal{L}= \underbrace{\mathrm{NMSE}_{\mathrm{DNA}}+\mathrm{NMSE}_{\mathrm{PROT}}}_{\text{reconstruction}}
+\lambda_{a}\,\|\bar z^{\mathrm{DNA}}_{s}-\bar z^{\mathrm{PROT}}_{s}\|^2
+\lambda_{c}\,\mathcal{L}_{\text{InfoNCE}}(z^{\mathrm{DNA}}_s,z^{\mathrm{PROT}}_s)
+\lambda_{o}\,\big(\mathrm{Cov}(z^{\mathrm{DNA}}_s,z^{\mathrm{DNA}}_p)+\mathrm{Cov}(z^{\mathrm{PROT}}_s,z^{\mathrm{PROT}}_p)\big),
\]
with $\bar z$ the $\ell_2$-normalized shared code. The InfoNCE term makes row $i$ of the DNA
shared code retrieve row $i$ of the protein shared code; because all variants are in BRCA1,
negatives are in-gene and retrieval cannot exploit gene identity. Orthogonality decorrelates
shared from private latents so that private codes capture modality-specific structure.

\paragraph{Causal intervention.} For shared latent $f$, its DNA decoder column
$D^{\mathrm{DNA}}_{s}[:,f]$ (rescaled to raw activation units) is a direction in Evo\,2's
residual stream; the protein column is a direction in ESM-2's. During a variant forward pass
we subtract the latent's reconstructed contribution at the variant position and re-score,
in each model, comparing against a matched-norm random direction.

\section{Experimental setup}
\label{sec:setup}
We use the BRCA1 SGE dataset \citep{findlay2018brca1}: \brcaNmis{} missense SNVs with
continuous DMS function scores, a functional class (FUNC/INT/LOF), and (partial) ClinVar
labels, mapped to UniProt P38398 (all \brcaNmis{} residue references validated). Variants
fall in two well-separated domains (RING, $n{=}\brcaRING{}$; BRCT, $n{=}\brcaBRCT{}$),
enabling a \emph{domain-disjoint} generalization split in addition to a residue-disjoint
split. Evo\,2 is \evoModel{} and ESM-2 is \esmModel{}; we select layers
$\ell_{\mathrm{DNA}}{=}\dnaLayer{}$ and $\ell_{\mathrm{PROT}}{=}\protLayer{}$ by a DMS probe
(\S\ref{sec:results}). The crosscoder uses shared/private width \Kshared{}/\Kprivate{} with
BatchTopK $k{=}\topk{}$. Baselines: single-model $\Delta$ linear probe, concatenation,
linear CCA, and PLS. All numbers are on held-out splits.

\section{Results}
\label{sec:results}
\paragraph{The alignment beats linear CCA; the sparse code is for interpretation.} Table~\ref{tab:main}. Our crosscoder's linear alignment head retrieves the paired variant at Recall@1 0.307$\pm$0.014 (5 seeds, residue-disjoint $n{=}\posNtest{}$), significantly above linear CCA (0.216; paired bootstrap $\Delta{=}0.104$, $p<0.001$) and matching a pure deep-CCA baseline (0.311). The \emph{sparse} shared code alone does not retrieve (0.003, chance) --- the cross-modal signal lives in the learned linear alignment, and the sparse codes serve interpretability (below). The same ordering holds for DMS (shared 0.442$\pm$0.003 $>$ CCA 0.337, DNA 0.375; ESM-2 0.514 and concat 0.514 remain stronger). Reconstruction FVE is 0.72 (DNA)/0.70 (protein).
\begin{table}[t]\centering\small
\caption{Cross-modal retrieval and DMS prediction (residue-disjoint, $n{=}\posNtest{}$; retrieval/DMS from 5 seeds). A learned alignment (deep-CCA or our crosscoder's alignment head) beats linear CCA at both; the crosscoder's \emph{sparse} shared code does not retrieve (it is the interpretable dictionary, not the alignment). ESM-2 alone / concat are the strongest DMS predictors.}\label{tab:main}
\begin{tabular}{lccc}
\toprule
Method & R@1 & R@10 & DMS $\rho$ \\
\midrule
linear CCA & 0.216 & 0.645 & 0.337 \\
deep CCA (align only) & 0.311 & 0.689 & 0.464 \\
\;crosscoder: sparse shared code & 0.003 & 0.027 & 0.272 \\
\;crosscoder: alignment head & 0.307 & 0.689 & 0.442 \\
ESM-2 $\Delta$ (probe) & -- & -- & 0.514 \\
concat / PLS & -- & -- & 0.514 \\
\bottomrule
\end{tabular}
\end{table}

\paragraph{Domain-disjoint generalization.} Training and evaluating only across the two disjoint domains (train BRCT, test RING; $n{=}\domNtest{}$), the alignment still beats CCA (Recall@1 0.106 vs.\ 0.081; DMS 0.303 vs.\ 0.279), though with substantial degradation (ESM-2 alone reaches 0.444). The method partially transfers to an unseen structural domain rather than memorizing residues.

\paragraph{Cross-gene generalization (multi-gene ClinVar).} To test the single-gene concern directly, we assemble \mgNgenes{} genes of ClinVar missense variants (balanced pathogenic/benign, GRCh38 coordinates, protein references validated) and train on \mgNtrainGenes{} genes, evaluating on \mgNtestGenes{} \emph{held-out} genes ($n{=}\mgNtest{}$). Cross-modal retrieval within each unseen gene reaches Recall@1 0.066 (Recall@10 0.340), above linear CCA (0.056) and comparable to deep-CCA (0.068); the alignment thus transfers to genes never seen in training. On \emph{gene-disjoint} ClinVar pathogenicity prediction, the shared code reaches AUROC 0.886, versus the Evo\,2 probe 0.864, the ESM-2 probe 0.855, and CCA 0.895---the shared cross-modal representation carries transferable pathogenicity signal across genes.

\paragraph{Biological structure of the sparse codes.} On held-out variants the shared sparse code predicts loss-of-function (AUROC 0.830, 95\% CI [0.776,0.877], $n{=}\lofN{}$) and domain (0.839, $n{=}\ringN{}$). Modality-private codes are complementary: the protein-private code is the best domain predictor (0.853) and matches shared on LOF (0.842), consistent with ESM-2 encoding structural context. The shared code also exceeds classical single-variant VEP predictors (SIFT 0.305, CADD 0.329, phyloP 0.336; $|\rho|$ vs.\ DMS).

\paragraph{A shared functional axis.} The direction of maximal DMS correlation, fit independently in each model's activation space, maps to nearly the same point in the alignment space (cosine 0.905), far above random direction pairs (0.306, 95th pctile $|{\cos}|=0.513$). This is property-specific: the two \emph{domain} directions also align (0.927), while a functional-vs-domain pair sits at the random level (0.282). Same biological property $\Rightarrow$ same shared axis across modalities; different properties $\Rightarrow$ different axes.

\paragraph{Causal probing (a negative result).} We test whether this alignment yields a causal handle by injecting the shared functional direction at the variant position and re-scoring ($n{=}\causalN{}$ held-out variants). Single-position activation edits do \emph{not} give reliable causal control in either model: the score shift is uncorrelated with the variant's true effect and indistinguishable from a matched-norm random direction (ESM-2 Spearman 0.131, $p{=}0.19$; Evo\,2 0.009, $p{=}0.93$). Representational alignment does not imply an interchangeable causal handle---a caution for cross-model steering, and an open question of whether richer interventions would succeed.

\section{Discussion and limitations}
Our study is deliberately scoped to one gene to eliminate gene-identity confounds; whether
shared latents transfer across genes and assays is the key next step (ProteinGym extension).
The paired-delta design assumes a well-defined variant position in both models, which holds
for missense SNVs but not for splice/indel variants, which we can only analyze through the
DNA branch. Causal effects are measured on the models' own likelihood scores, not on
downstream phenotype.

\section{Conclusion}
Paired variant deltas plus a shared--private crosscoder expose sparse, causally active
mechanisms shared between a genomic and a protein foundation model, and separate them from
modality-specific ones---a step toward mechanistic, cross-modal interpretability of biological
foundation models.

\small
\bibliographystyle{plainnat}
\bibliography{refs}
\end{document}
